\documentclass[9pt,a4paper]{ctexart}
\usepackage[margin=9mm]{geometry}
\usepackage{enumitem,hyperref,tabularx,xcolor,titlesec}
\usepackage{fontspec}

\setmainfont{Helvetica Neue}
\setCJKmainfont{STHeiti}
\hypersetup{colorlinks=true,urlcolor=blue!65!black}
\pagestyle{empty}
\setlength{\parindent}{0pt}
\setlength{\parskip}{0pt}
\setlist[itemize]{leftmargin=1.2em,itemsep=0pt,topsep=1pt,parsep=0pt,partopsep=0pt}
\titleformat{\section}{\bfseries\fontsize{13}{14}\selectfont}{}{0pt}{}[\vspace{-2pt}\titlerule]
\titlespacing*{\section}{0pt}{5pt}{3pt}
\newcommand{\period}[1]{\hfill\textcolor{gray}{\small #1}}
\newcommand{\tags}[1]{\textcolor{black!75}{\footnotesize #1}}
\newcommand{\labelitem}[2]{\item \textbf{#1：}#2}

\begin{document}
\fontsize{9.2}{11.2}\selectfont

\begin{tabularx}{\textwidth}{@{}Xr@{}}
  {\fontsize{22}{24}\selectfont\bfseries CHZarles} &
  \begin{tabular}{r}
    \href{mailto:2538745263@qq.com}{2538745263@qq.com}\\
    \href{tel:15697580245}{15697580245}\\
    \href{https://github.com/CHZarles}{github.com/CHZarles}
  \end{tabular}
  \\
  \textcolor{black!70}{中国} &
\end{tabularx}

\section*{工作经历}

\textbf{瀚博半导体}\period{2022/07 -- 2025/07}\\[-1pt]
\textbf{项目一：自研 NPU 高性能推理 SDK 与 Python 绑定开发}\\[-1pt]
\tags{C++ \quad CMake \quad Python \quad pybind11 \quad NumPy Buffer Protocol \quad moodycamel \quad libgo \quad RAII}

面向公司自研 NPU 芯片的推理接入需求，针对底层 C Runtime 接口原始、异步回调侵入性强及手动管理 Buffer 易泄漏等问题，基于 CMake 设计并实现跨语言 C++/Python 推理 SDK。将底层异步回调转换为拉取式结果接口，并提供与 NumPy 互通的 Pythonic API，为算法与应用团队提供高效、内存安全的推理开发套件。
\begin{itemize}
  \labelitem{异步结果通道与拉取式接口}{接管 C Runtime completion callback，以 moodycamel 无锁并发队列作为跨线程结果通道，并通过 ProducerToken 优化回调路径；业务线程仅需调用 RunAsync() 提交任务，通过 GetOutput(timeout) 阻塞等待并拉取完整 batch 结果。}
  \labelitem{回调解耦与异步调度}{回调函数仅负责组装 Result 并高效入队，与业务结果处理严格解耦；SDK 内部结合 libgo 承载异步调度基础设施，避免底层回调逻辑侵入上层业务代码。}
  \labelitem{Buffer 生命周期安全闭环}{使用 RAII 与智能指针封装模型句柄、执行流上下文及 Tensor/Image Buffer，任务对象自动持有执行期间所需资源；设计 CloseInput() \(\rightarrow\) 持续 GetOutput() 至 EOF \(\rightarrow\) WaitUntilDone() 的收尾状态机，待 Runtime 回调静止、消费端退出后再销毁队列与上下文，规避悬垂指针和内存泄漏。}
  \labelitem{Pythonic 跨语言封装}{基于 pybind11 封装 C++ 类、上下文与推理生命周期，提供上下文管理器 with、流式迭代和异常映射，为上层应用屏蔽底层并发与内存管理细节。}
  \labelitem{NumPy 互通与并发安全}{基于 Python Buffer Protocol 实现 Tensor Buffer 与 numpy.ndarray 的零拷贝/低开销互通；在 GetOutput() 阻塞路径显式释放 GIL（py::gil\_scoped\_release），避免阻塞 Python 主线程及并发业务逻辑。}
  \labelitem{商业级交付验证}{SDK 作为一体化解决方案核心组件，交付至阿里云、小红书、高视科技等客户，并在内部算法压测及生产业务中持续运行。}
\end{itemize}

\vspace{2pt}
\textbf{瀚博半导体}\period{2022/07 -- 2025/07}\\[-1pt]
\textbf{项目二：端到端 AI 业务管线构建}\\[-1pt]
\tags{Python \quad C++ \quad MediaPipe \quad OpenCV \quad PyTorch \quad YOLOv8}
\begin{itemize}
  \labelitem{端到端算法流程落地}{负责目标检测（YOLOv8）、特征提取（ResNet）及多目标跟踪（ByteTrack）等算法在自研硬件上的全流程打通；基于 OpenCV、NumPy、LibTorch 编写高性能前后处理及 NPU 协同推理代码。}
  \labelitem{复杂业务流编排}{深入 Google MediaPipe，根据业务需求开发多种 Calculator 节点，实现数据输入、前后处理、多模型串联推理及跟踪逻辑的计算图编排。}
  \labelitem{环境容器化与部署标准化}{构建 Docker 镜像并收敛依赖，整合 NPU 运行时及 C++/Python 依赖；规范模型仓库挂载与数据卷映射，确保不同环境下一致、快速部署。}
  \labelitem{工程文档与交付支持}{编写端到端调用范例与接入文档，协助业务团队及外部客户快速完成算法接入和问题排查。}
\end{itemize}

\section*{项目经历}

\textbf{PaperLoom 论文智能阅读与可溯源问答系统}
\quad\href{https://paperloom.me/}{paperloom.me}\\[-1pt]
\tags{OpenAI Agents SDK \quad Java 17 / Spring Boot \quad Qdrant \quad MySQL \quad Kafka \quad MinIO \quad WebSocket}

面向科研论文阅读场景开发 Agentic RAG 问答系统。采用 Java/Spring Boot（业务边界与数据权威）+ Python（Agent 运行时）分层架构，确保回答中的每一句论述均可回溯至原始 PDF 的页码、章节与段落。
\begin{itemize}
  \labelitem{Agent 工具链与引用防伪校验}{设计“候选检索 \(\rightarrow\) 原文精读 \(\rightarrow\) 引用生成 \(\rightarrow\) 最终提交”调用链；维护已读原文清单并以代码规则校验引用编号，从机制上杜绝虚假引用。}
  \labelitem{数据分层与结构化建模}{对 PDF 解析产物进行领域建模，统一为带页码、章节层级、表格/图表及坐标的数据结构；以 Qdrant 导航候选位置、MySQL 提供权威原文。}
  \labelitem{权限锁定与异步解析流水线}{以数据库锁固定会话论文范围；基于 Kafka、MinIO、MySQL 状态机构建分片上传、异步解析与向量索引流水线，通过唯一约束和条件更新保障幂等。}
  \labelitem{流式交互与离线评测}{基于 Spring WebSocket 流式推送 Agent 思考过程、工具状态与答案，支持任务中断；搭建 Golden Cases，分层统计检索召回率、阅读覆盖率和引用准确率。}
\end{itemize}

\section*{教育经历}

\begin{tabularx}{\textwidth}{@{}Xr@{}}
  \textbf{澳门大学} \quad 人工智能 · 硕士 & \textcolor{gray}{2025/08 -- 2027/06}\\[2pt]
  \textbf{中国矿业大学} \quad 计算机科学与技术 · 本科，GPA: 87/100，专业排名：26/139 & \textcolor{gray}{2018/09 -- 2022/07}
\end{tabularx}

\end{document}
